\begin{abstract}
A question about a new prime minister, CEO or coach can contain a true premise that conflicts with a language model's outdated knowledge. We study how seven small open models (0.5--7B parameters, released in 2024) respond to such newly-true presuppositions. Our audited Wikidata benchmark includes matched triples: the current holder appointed since 2025 (true), a former holder (false), and another entity's new holder swapped in (false), all asked about the same role with identical wording. Without added instructions, models rebut \WGNewPlain\% of premises about recent changes and \WGOldPlain\% of long-true ones. Within triples, they rebut true premises about ten percentage points less often than false ones (\WNPlain\% vs.\ \WAPlain\% and \WSPlain\%). This margin is positive in five models and zero in two; the evidence for discrimination is suggestive. Most rebuttals invoke a knowledge-cutoff hedge; unhedged rebuttals of true premises occur in \WUhN\% of answers. Instruction wording matters: our premise check raises rebuttals of true and false premises alike, whereas FreshQA's sentence leaves the true-premise rate nearly unchanged and widens the margin to the swapped holder in six models. Stating the cutoff shifts answers toward deferral; an oracle snippet naming the current holder removes false rebuttals. Three larger models (14--32B) rebut recent changes about as often as the small models that show prior role knowledge, but rebut long-true premises less often. We check the primary model judge against held-out labels from Claude Sonnet~5.5 and OpenAI gpt-5.3-codex, and release a regenerable benchmark.
\end{abstract}
